\documentclass[12pt,a4paper]{article}

\usepackage[utf8]{inputenc}
\usepackage[T1]{fontenc}
\usepackage{lmodern}
\usepackage[spanish]{babel}
\usepackage{geometry}
\usepackage{enumitem}
\usepackage{xcolor}
\usepackage[hidelinks]{hyperref}

\input{glyphtounicode}
\pdfgentounicode=1

\geometry{margin=2.5cm}
\setlist[itemize]{leftmargin=1.2cm,itemsep=0.25em}
\setlist[enumerate]{leftmargin=1.2cm,itemsep=0.25em}
\definecolor{imegreen}{RGB}{15,118,110}
\newcommand{\dato}[2]{\noindent\textbf{#1:} #2\par\vspace{0.25em}}

\title{\textbf{Bases Oficiales}\\[0.35cm]
Campeonato Deportivo de Ingeniería Mecánica Eléctrica 2026}
\author{Comisión Deportiva de Ingeniería Mecánica Eléctrica}
\date{Actualizado al 1 de julio de 2026}

\begin{document}

\pagenumbering{gobble}
\begin{titlepage}
  \centering
  \vspace*{1.5cm}
  {\Large \textbf{Escuela Profesional de Ingeniería Mecánica Eléctrica}\par}
  \vspace{1.2cm}
  {\Huge \textbf{Bases Oficiales del Campeonato}\par}
  \vspace{0.35cm}
  {\LARGE Fútbol 11 y Vóley Mixto\par}
  \vspace{1.2cm}
  \rule{0.72\textwidth}{1.2pt}
  \vspace{0.9cm}

  \begin{minipage}{0.82\textwidth}
    \dato{Nombre del campeonato}{Campeonato Deportivo de Ingeniería Mecánica Eléctrica 2026}
    \dato{Organiza}{Comisión Deportiva de Ingeniería Mecánica Eléctrica}
    \dato{Lugar}{Campus universitario y canchas designadas por la organización}
    \dato{Costo de inscripción}{Fútbol 11: S/ 35.00 por equipo. Vóley Mixto: S/ 25.00 por equipo}
    \dato{Medios de pago}{Yape o Plin, mediante código único de inscripción}
    \dato{Contacto}{Comisión Deportiva de Ingeniería Mecánica Eléctrica / [celular o correo oficial]}
  \end{minipage}

  \vfill
  {\large Documento listo para impresión, descarga digital y carga en la sección de Bases del sistema.\par}
  \vspace{1cm}
\end{titlepage}

\pagenumbering{roman}
\tableofcontents
\newpage
\pagenumbering{arabic}

\section{Presentación}

Las presentes bases regulan la organización y participación en el Campeonato Deportivo de Ingeniería Mecánica Eléctrica 2026, que comprende las disciplinas de Fútbol 11 y Vóley Mixto. El campeonato busca promover la integración, el respeto, la competencia sana y la participación responsable de los equipos inscritos.

La inscripción, programación, registro de jugadores, resultados, sanciones y publicación de información oficial se administrarán mediante el sistema del campeonato y por las decisiones de la organización.

\section{Objetivos}

\begin{enumerate}
  \item Fomentar la práctica deportiva entre estudiantes y equipos vinculados a Ingeniería Mecánica Eléctrica.
  \item Promover la integración, disciplina, puntualidad y juego limpio.
  \item Establecer reglas claras para inscripción, participación, arbitraje y resolución de reclamos.
  \item Publicar resultados y avances del campeonato de forma transparente.
\end{enumerate}

\section{Datos generales del campeonato}

\dato{Institución / escuela}{Escuela Profesional de Ingeniería Mecánica Eléctrica}
\dato{Campeonato}{Campeonato Deportivo de Ingeniería Mecánica Eléctrica 2026}
\dato{Deportes}{Fútbol 11 y Vóley Mixto}
\dato{Lugar}{Campus universitario y canchas designadas por la organización}
\dato{Costo de inscripción}{Fútbol 11: S/ 35.00 por equipo. Vóley Mixto: S/ 25.00 por equipo}
\dato{Pago}{Yape o Plin. El encargado entrega un código único de inscripción después de verificar el pago.}
\dato{Cierre de inscripción Fútbol 11}{8 de julio de 2026, 23:59 horas, salvo modificación comunicada por la organización.}
\dato{Cierre de inscripción Vóley Mixto}{10 de julio de 2026, 23:59 horas, salvo modificación comunicada por la organización.}
\dato{Contacto oficial}{Comisión Deportiva de Ingeniería Mecánica Eléctrica / [celular o correo oficial]}
\dato{Premios}{[Premios oficiales a confirmar por la organización]}
\dato{Regla mixta definitiva de vóley}{[Regla de composición mixta pendiente de confirmación por la organización]}

\section{Inscripción general}

\subsection{Condiciones para inscribirse}

Solo podrán participar equipos inscritos dentro del plazo establecido y con pago validado por la organización. El registro se realiza en el sistema oficial del campeonato, utilizando el código único entregado luego del pago.

\subsection{Costo y pago}

\begin{itemize}
  \item El costo de inscripción para Fútbol 11 es de \textbf{S/ 35.00 por equipo}.
  \item El costo de inscripción para Vóley Mixto es de \textbf{S/ 25.00 por equipo}.
  \item El costo de arbitraje por partido es de \textbf{S/ 20.00} para Fútbol 11 y \textbf{S/ 10.00} para Vóley Mixto.
  \item El pago se realiza por \textbf{Yape} o \textbf{Plin}, según indique la organización.
  \item Luego de verificar el pago, la organización entregará un \textbf{código único de inscripción}.
  \item El código solo puede usarse una vez y queda asociado al equipo inscrito.
  \item La inscripción está sujeta a validación de pago y revisión de requisitos.
\end{itemize}

\subsection{Datos obligatorios del equipo}

Cada equipo debe registrar:

\begin{itemize}
  \item nombre del equipo;
  \item deporte y categoría;
  \item nombre completo del delegado;
  \item número de celular del delegado;
  \item correo del delegado;
  \item carrera, escuela o grupo al que representa, si corresponde;
  \item método de pago utilizado;
  \item código único de inscripción.
\end{itemize}

\subsection{Datos obligatorios de cada jugador}

Cada jugador debe registrar:

\begin{itemize}
  \item nombres;
  \item apellidos;
  \item DNI;
  \item código de estudiante o jugador;
  \item semestre o ciclo académico;
  \item ficha de matrícula en formato PDF, JPG o PNG;
  \item número de camiseta, si la organización lo solicita;
  \item rol o condición deportiva, si corresponde.
\end{itemize}

\subsection{Archivo de ficha de matrícula}

La ficha de matrícula es obligatoria para cada jugador. La organización podrá rechazar u observar registros sin ficha, con ficha ilegible o con información inconsistente.

El sistema admitirá archivos PDF, JPG o PNG. El tamaño máximo recomendado por archivo es de 5 MB.

\section{Requisitos de equipos}

\begin{itemize}
  \item Cada equipo debe tener un delegado responsable.
  \item Los equipos deben completar la cantidad mínima de jugadores exigida por su deporte.
  \item Solo se considerarán habilitados los equipos con pago validado y aprobación de la organización.
  \item La organización podrá observar, rechazar o descalificar inscripciones incompletas o falsas.
  \item No se permite suplantación de jugadores ni documentación adulterada.
\end{itemize}

\section{Requisitos de jugadores}

\begin{itemize}
  \item El jugador debe estar correctamente registrado en el sistema.
  \item No se permite que un jugador participe en dos equipos de la misma categoría.
  \item La duplicidad se verificará por DNI y por código de estudiante o jugador.
  \item El jugador debe presentar datos verdaderos y ficha de matrícula válida.
  \item Un jugador no habilitado puede ocasionar observación, pérdida del partido o descalificación del equipo, según la gravedad.
\end{itemize}

\section{Rol del delegado}

El delegado es el representante oficial del equipo ante la organización. Sus responsabilidades son:

\begin{itemize}
  \item registrar correctamente al equipo y sus jugadores;
  \item revisar que cada jugador tenga DNI, código, semestre y ficha de matrícula;
  \item recibir y atender comunicaciones oficiales;
  \item presentar reclamos, si correspondiera;
  \item garantizar puntualidad y conducta deportiva de su equipo.
\end{itemize}

Mientras la inscripción esté abierta y el campeonato no haya iniciado, el delegado podrá solicitar o realizar modificaciones permitidas por el sistema. Una vez iniciado el campeonato, el delegado no podrá editar datos del equipo ni del plantel.

Después del inicio, solo podrá modificar el número de camiseta de cada jugador una única vez, si el sistema y la organización lo permiten. El número debe estar entre 1 y 99 y no puede repetirse dentro del mismo equipo.

\section{Campeonato de Fútbol 11}

\subsection{Datos específicos}

\dato{Deporte}{Fútbol 11}
\dato{Formato}{Eliminación directa}
\dato{Costo de inscripción}{S/ 35.00 por equipo}
\dato{Fecha de inicio}{9 de julio de 2026}
\dato{Hora de inicio}{8:00 a. m.}
\dato{Canchas}{1 cancha}
\dato{Fixture}{Programación o sorteo aleatorio realizado por la organización}
\dato{Equipos participantes}{Solo equipos inscritos, con pago validado y aprobados por la organización}

\subsection{Duración del partido}

Cada partido tendrá:

\begin{itemize}
  \item primer tiempo: 10 minutos;
  \item descanso: 10 minutos;
  \item segundo tiempo: 15 minutos.
\end{itemize}

El árbitro será responsable del control del tiempo oficial y de registrar el estado del partido en el sistema.

\subsection{Empate}

Al tratarse de eliminación directa:

\begin{itemize}
  \item si el partido termina empatado en el tiempo reglamentario, se definirá por penales;
  \item no habrá tiempo extra, salvo decisión expresa de la organización;
  \item el ganador por penales avanzará a la siguiente fase;
  \item si se trata de la final, el ganador será declarado campeón.
\end{itemize}

\subsection{Penales}

\begin{itemize}
  \item Cada equipo ejecutará una tanda inicial de 3 penales.
  \item Si persiste el empate, se continuará con muerte súbita.
  \item Los goles de penales no se suman al marcador del tiempo reglamentario.
  \item Los penales solo sirven para definir al ganador.
\end{itemize}

Ejemplo de resultado:

\begin{center}
  \fbox{\textbf{Equipo A 1 (3) -- 1 (2) Equipo B}}
\end{center}

\subsection{Arbitraje en Fútbol 11}

\begin{itemize}
  \item El árbitro es la autoridad principal del partido.
  \item El costo de arbitraje por partido es de \textbf{S/ 20.00}.
  \item El resultado cargado por el árbitro se publicará inmediatamente como resultado oficial.
  \item La organización podrá revisar o corregir el resultado solo en caso de controversia, reclamo o error comprobado.
  \item Las tarjetas rojas generan suspensión automática mínima de un partido.
  \item Las tarjetas amarillas se registrarán para control disciplinario.
\end{itemize}

\section{Campeonato de Vóley Mixto}

\subsection{Datos específicos}

\dato{Deporte}{Vóley Mixto}
\dato{Modalidad}{Mixto}
\dato{Estado}{Abierto a inscripción}
\dato{Sistema de juego}{Por sets}
\dato{Formato}{Editable por la organización según cantidad de equipos inscritos}
\dato{Costo de inscripción}{S/ 25.00 por equipo}

\subsection{Puntuación de sets}

Para que el campeonato sea rápido, se aplicará la siguiente puntuación:

\begin{itemize}
  \item set 1: a 15 puntos;
  \item set 2: a 15 puntos;
  \item set 3: a 10 puntos, solo si hay empate 1--1.
\end{itemize}

Regla de diferencia:

\begin{itemize}
  \item cada set debe ganarse por diferencia mínima de 2 puntos;
  \item set 1 y set 2 tendrán tope máximo de 17 puntos;
  \item set 3 tendrá tope máximo de 12 puntos.
\end{itemize}

Ejemplos:

\begin{itemize}
  \item si un set va 14--14, se juega hasta diferencia de 2;
  \item si llega a 16--16, gana quien haga el punto 17;
  \item en el tercer set, si llega a 11--11, gana quien haga el punto 12.
\end{itemize}

\subsection{Composición mixta}

Cada equipo deberá mantener al menos [número] jugadores/as de cada género en cancha, según defina la organización.

Mientras no exista definición final, se aplicará el siguiente campo editable:

\begin{center}
  \fbox{\parbox{0.86\textwidth}{\centering
  [Regla de composición mixta pendiente de confirmación por la organización]}}
\end{center}

\subsection{Resultados y Arbitraje de Vóley Mixto}

\begin{itemize}
  \item El árbitro o responsable registrará los sets ganados y los puntos por set.
  \item El costo de arbitraje por partido es de \textbf{S/ 10.00}.
  \item El resultado se publicará en el sistema una vez registrado.
  \item La organización podrá revisarlo en caso de controversia.
\end{itemize}

\section{Arbitraje y resultados}

\begin{itemize}
  \item El árbitro o responsable designado controlará el desarrollo del partido.
  \item El costo de arbitraje por partido es de \textbf{S/ 20.00} para Fútbol 11 y \textbf{S/ 10.00} para Vóley Mixto.
  \item Los resultados registrados por el árbitro se publicarán como resultado oficial.
  \item El público, delegados, árbitros y administradores podrán visualizar los resultados publicados.
  \item La organización solo intervendrá después si existe controversia, reclamo, error reportado u observación administrativa.
  \item Las correcciones deberán quedar registradas con motivo, fecha y responsable.
\end{itemize}

\section{Sanciones}

\begin{itemize}
  \item Tarjeta roja directa: suspensión mínima de 1 partido.
  \item Agresión verbal o física: evaluación por la organización y posible suspensión mayor o descalificación.
  \item Jugador no habilitado: posible pérdida del partido.
  \item Inscripción falsa: descalificación del equipo.
  \item Conducta antideportiva: sanción según gravedad.
  \item La organización podrá ampliar sanciones en casos de reincidencia o hechos graves.
\end{itemize}

\section{Reclamos}

\begin{itemize}
  \item Los reclamos deberán ser presentados únicamente por el delegado.
  \item Deben realizarse dentro de un plazo máximo de [tiempo] después del partido.
  \item La organización revisará evidencias, reporte arbitral y registros del sistema.
  \item La decisión de la organización será final.
\end{itemize}

\section{Premiación}

\begin{itemize}
  \item Campeón Fútbol 11: [premio]
  \item Subcampeón Fútbol 11: [premio]
  \item Goleador Fútbol 11: [premio]
  \item Equipo menos goleado Fútbol 11: [premio]
  \item Campeón Vóley Mixto: [premio]
  \item Subcampeón Vóley Mixto: [premio]
  \item Mejor equipo Vóley Mixto: [premio]
\end{itemize}

\section{Disposiciones finales}

\begin{enumerate}
  \item La participación en el campeonato implica la aceptación total de estas bases.
  \item Los casos no previstos serán resueltos por la organización.
  \item La organización podrá ajustar horarios, canchas o fixture por razones justificadas.
  \item Cualquier modificación será comunicada por los canales oficiales.
  \item La organización velará por la transparencia, seguridad y correcto desarrollo del campeonato.
\end{enumerate}

\vfill
\begin{center}
  \textcolor{imegreen}{\textbf{Comisión Deportiva de Ingeniería Mecánica Eléctrica}}\\
  Campeonato Deportivo de Ingeniería Mecánica Eléctrica 2026
\end{center}

\end{document}
